\section{From simple-curve coefficients to controlled observables}
\label{sec:coefficientobservables}

The preceding section supplies an explicit analytic norm for its mixing
estimate. We now connect that norm to a selected family of simple curves
and their coefficient distributions. The genus-one norm-ball obstruction
first fixes the attainable scale of sharp remainders. Primitive
pole-to-pole arcs then give geometric coefficient observables with
uniform rotation and terminal-derivative bounds. Exact affine Gaussian
calculations provide independent tests of the analytic estimates.
Throughout, a fixed-surface sharp count, a moduli-averaged covariance,
and the conditioning of coefficient extraction remain separate problems.

\subsection{The norm-ball obstruction and the correct error scale}
\label{subsec:normballbarrier}

The genus-one length norm gives a necessary benchmark for any passage from
smooth return estimates to sharp simple-curve counting. Let $X$ be a fixed
complete hyperbolic once-punctured torus, and let $B_X$ be the unit ball of
the McShane--Rivin length norm on $H_1(X;\RR)\simeq\RR^2$. Put
\[
 A_X(L)=\#(LB_X\cap\ZZ^2),\qquad
 P_X(L)=\#(LB_X\cap\ZZ^2_{\mathrm{prim}}).
\]
The first count includes the origin and nonprimitive classes. The second
counts oriented primitive classes; the count of unoriented simple closed
geodesics is $N_X(L)=\tfrac12P_X(L)$. Elementary planar lattice counting
and M\"obius inversion give
\begin{align}
 A_X(L)&=\operatorname{Area}(B_X)L^2+O_X(L),\label{eq:normballall}\\
 P_X(L)&=\frac{\operatorname{Area}(B_X)}{\zeta(2)}L^2
           +O_X(L\log L).\label{eq:normballprimitive}
\end{align}
Indeed, $P_X(L)=\sum_d\mu(d)(A_X(L/d)-1)$, with $d\ll_X L$;
the boundary errors sum as $L\sum_{d\ll_X L}d^{-1}$. These are the
upper estimates behind \cite[Theorem~4.1]{McShaneRivinTori}, with the
orientation factor made explicit here.

The next theorem in that paper concerns a dense set of hyperbolic
structures. Its mechanism is an infinitely flat boundary point whose
supporting normal is rational with respect to the homology lattice.
The subsequent Conjecture~4.3 proposes an $O(L^{1/2+\varepsilon})$ error
for the modular torus; thus the obstruction is not asserted there for
every individual surface. The phrase ``at least $O(L)$'' in Theorem~4.2
should be separated from the literal modern assertion
$\limsup|E_X(L)|/L>0$. The elementary argument below proves the precise
consequence needed here: no fixed improvement of the exponent below one
is possible under the stated flat-point hypothesis. We do not use it to
claim a literal $\Omega(L)$ bound.

\begin{proposition}[A primitive-row obstruction]
\label{prop:normballflatrow}
Let $B\subset\RR^2$ be a strictly convex compact body with the origin
in its interior. Suppose that $p\in\partial B$ has irrational radial
direction, a rational supporting normal, and infinite graph-flatness.
Then, for
\[
 E_B(L)=\#(LB\cap\ZZ^2_{\mathrm{prim}})
       -\frac{\operatorname{Area}(B)}{\zeta(2)}L^2,
\]
and every $\varepsilon>0$, there are $c_\varepsilon>0$ and arbitrarily
large $L$ such that
\begin{equation}
 |E_B(L)|\geq c_\varepsilon L^{1-\varepsilon}.
 \label{eq:normballlowerpower}
\end{equation}
In particular $E_B(L)=O(L^{1-\eta})$ fails for every fixed $\eta>0$.
\end{proposition}

\begin{proof}
An integral unimodular change of coordinates makes the supporting normal
vertical. Write $p=(\alpha,h)$, with $h>0$ and $\alpha/h$ irrational.
Near $p$, the upper boundary has the form
\[
 y=h-f(x-\alpha),\qquad f(0)=0,
\]
where $f\geq0$ is convex. Infinite graph-flatness means that for every
integer $N>2$ there are $C_N,\rho_N>0$ such that
\[
 0\leq f(s)\leq C_N|s|^N\qquad (|s|\leq\rho_N).
\]
Fix $N$, let $n$ run through primes, and set
\[
 L_n=\frac nh,\qquad \delta_n=L_n^{-1},\qquad L_n^+=L_n+\delta_n.
\]
The line $y=n$ supports $L_nB$ at the single point
$(L_n\alpha,n)$. Its first coordinate is irrational, so this row
contains no lattice point of $L_nB$.

In $L_n^+B$, the same row lies a normalized depth
$h\delta_n/L_n^+$ below the supporting line. Therefore it contains the
horizontal interval centered at $L_n^+\alpha$ with half-width
\[
 r_n=\min\left\{
 \left(\frac h{2C_N}\right)^{1/N}
 (L_n^+)^{1-1/N}\delta_n^{1/N},\ \frac n4
 \right\}.
\]
For sufficiently large $n$ the local graph description applies, the
first quantity realizes the minimum, and
\[
 r_n\asymp_{B,N}L_n^{1-2/N}.
\]
The interval has length at most $n/2$. Among its integer points
$(k,n)$, at most one has $n\mid k$; every other point is primitive
because $n$ is prime. Thus
\[
 P_B(L_n^+)-P_B(L_n)\geq 2r_n-O(1).
\]
The change in the quadratic main term is bounded, since
\[
 (L_n^+)^2-L_n^2=2+L_n^{-2}.
\]
It follows that
\[
 \max\{|E_B(L_n)|,|E_B(L_n^+)|\}
 \geq c_{B,N}L_n^{1-2/N}
\]
for all sufficiently large primes $n$. Choose $N>2/\varepsilon$ to
obtain the assertion. To rule out a proposed $O(L^{1-\eta})$ estimate,
apply the assertion with $0<\varepsilon<\eta$.
\end{proof}

The proof works directly with primitive classes, so no cancellation in a
M\"obius sum needs to be controlled. The short interval
$[L_n,L_n+L_n^{-1}]$ acquires many new primitive points along a nearly
flat supporting row, while its predicted main-term increment stays
bounded. Infinite flatness alone does not force the row length to be a
positive fraction of $L_n$; that is why the proof gives an exponent-one
barrier, not a literal linear lower bound. The distinction between a
fixed-surface estimate, an estimate uniform over moduli, and a smoothed
estimate must be retained.

\subsection{The Fourier and spectral meaning of the obstruction}
\label{subsec:normballspectral}

The flatness obstruction has a direct oscillatory-integral interpretation.
For illustration, suppose a localized smooth boundary arc has graph
\[
 y=h-f(x),\qquad f(x)=a x^k+O(x^{k+1}),\quad a>0,
\]
where $k\geq2$ is even. Apply smooth cutoffs supported near this arc
and equal to one at the supporting point. Integrating the localized
indicator in the normal coordinate first gives, up to rapidly decreasing
terms,
\[
 -\frac{e^{-2\pi i r h}}{2\pi i r}
       \int \chi(x)e^{2\pi i r f(x)}\,\dd x.
\]
The substitution $x=r^{-1/k}u$ in the stationary-phase calculation gives
\begin{equation}
 \widehat{\chi_0\mathbf1_B}(r\nu)
   =c_{a,k}\,e^{-2\pi i r h}r^{-1-1/k}
     +o(r^{-1-1/k}),\qquad c_{a,k}\ne0.
 \label{eq:normballfinitephase}
\end{equation}
where $\nu$ is the chosen normal direction and $\chi_0$ is the localized
cutoff. The familiar nondegenerate exponent is $r^{-3/2}$ when $k=2$;
larger contact order weakens it toward $r^{-1}$. Localization matters:
the contributions of different supporting points can interfere in the
Fourier transform of the whole ball. No global lower bound follows from
discarding those phases.

Under dilation,
\[
 \widehat{\mathbf1_{LB}}(m)=L^2\widehat{\mathbf1_B}(Lm).
\]
Thus a finite-type flat point with a lattice-rational normal contributes
on the scale $L^{1-1/k}$ to the corresponding localized Poisson terms.
For an infinitely flat point there is no fixed finite type to insert in
this formula. Proposition~\ref{prop:normballflatrow} supplies a direct
counting argument in that case, without assuming smooth boundary
regularity at all nearby points or invoking a formal stationary-phase
expansion there.

There is also an exact elementary spectral model. On an auxiliary flat
torus $\RR^2/\ZZ^2$, define the positive self-adjoint Fourier multiplier
\[
 \mathsf P_X e^{2\pi i\langle m,x\rangle}
   =\|m\|_X e^{2\pi i\langle m,x\rangle}.
\]
Norm equivalence gives a discrete spectrum tending to infinity. Its
counting function is $A_X(L)$, and restriction to the closed span of
primitive Fourier modes has counting function $P_X(L)$. Restricting
further to the even part identifies $m$ and $-m$ and gives $N_X(L)$.
This multiplier is defined spectrally; the nonsmooth length norm is not
being asserted to be a classical smooth pseudodifferential symbol.
At a differentiable point $p$ of its unit ball, a supporting covector
$\nu$ satisfies
\[
 \nabla\|\cdot\|_X(p)=\frac{\nu}{\langle\nu,p\rangle}.
\]
Thus a lattice-rational normal is exactly a rational direction for the
translation velocity of the associated Hamiltonian. This explains the
periodic-phase interpretation, without assuming that a global smooth
trace formula applies to the length norm.

The norm-ball discrepancy is therefore a spectral-projector remainder
for this particular length multiplier. Its flat normal directions
produce degenerate phases in a regularized Poisson or trace calculation.
This is distinct from complementary-series spectrum for the
$SL_2(\RR)$ action on the marked-torus moduli space. In particular, it
does not contradict the temperedness statement or the smooth
correlation estimate in Section~\ref{sec:genus2markedtorus}. The
rotation-Sobolev issue depends on what is held fixed. A sharp indicator
of a nonradial region in fixed physical vector coordinates generally
has rotational jumps, so smoothing it introduces growing derivative
norms. By contrast, hyperbolic length is unchanged when the flat charts
are rotated while the underlying conformal surface stays fixed. A sharp
cutoff in that invariant length can therefore lie in the rotation
generator's domain, and a radial flat cutoff can be rotation invariant
as well. Sharpness alone is not an obstruction to the stated mixing
bound. For those invariant observables, the remaining boundary and
regularity issues enter when taking flow derivatives or passing from
an averaged correlation to a fixed-surface counting remainder.

\subsection{A coefficient observable from primitive pole-to-pole arcs}
\label{subsec:primitivepolebridge}

The marked-torus construction supplies an explicit family of simple curves
for which the coefficient data become observables with controlled norms.
This family carries a flat-length cutoff as well as a hyperbolic-length
weight. The two length parameters remain distinct.

Work on the finite labeled model $Y_2$ of
Section~\ref{subsec:genus2homogeneous}. Write
$E=\RR^2/g\ZZ^2$, with its origin one of the four ramification points
of the pillowcase map. Fix a nonzero parity vector
$\delta\in(\ZZ/2\ZZ)^2$. For a nonzero primitive integer vector
$m\equiv\delta\pmod2$, project the straight segment
\[
 [0,gm/2]\subset\RR^2
\]
first to $E$ and then to $E/\{\pm1\}$. It joins the pole at the
origin to the pole corresponding to $\delta/2$.

\begin{lemma}[Primitive directions give simple curves]
\label{lem:primitivepolearcs}
Outside a Haar-null invariant subset of $Y_2$, the preceding segment
projects to an embedded arc avoiding the two regular marked points.
Its lift to the genus-two double cover is a nonseparating simple
closed regular flat geodesic $\alpha_m$, with
\begin{equation}
 \ell_{q_X}(\alpha_m)=|gm|.
 \label{eq:polearcflatlength}
\end{equation}
Opposite vectors give the same curve; distinct primitive directions
modulo sign give distinct free homotopy classes.
\end{lemma}

\begin{proof}
If two parameters $0<s,t<1$ on the segment identify on the pillowcase,
then $(s-t)m\in2\ZZ^2$ or $(s+t)m\in2\ZZ^2$. Primitivity forces
the corresponding scalar to be an even integer. The first case gives
$s=t$, and the second is impossible. Likewise the interior contains no
two-torsion point. The segment hits one of the additional regular marks
only if one of its lifts $p,q$ lies on the closed rational circle
$\RR gm/g\ZZ^2$. For each $m$ this is a null condition in the point
fiber. The union over the countable set of primitive $m$ is still null
and is invariant under the entire $SL_2(\RR)$ action.

The preimage of an embedded arc joining two branch values of a double
cover is a circle. Removing it leaves a connected cover, since the four
other branch values still have nontrivial monodromy; thus the circle is
nonseparating. At its endpoints the sphere flat angle $\pi$ doubles to
$2\pi$, so the two lifted segments join to a regular flat geodesic.
Its length is twice $|gm|/2$. Different directions give intersecting
regular geodesics with distinct tangents at the Weierstrass point over
the origin. Freely homotopic regular geodesics in a nonpositively curved
flat metric lie in a flat cylinder and cannot intersect transversely.
This proves the last assertion.
\end{proof}

Write $\ell_m(q_X)=\ell_X(\alpha_m)$ for the hyperbolic length of the
same topological curve. Proposition~\ref{prop:coefffirstatom} identifies
it as the position of the first atom of $\sigma^X_{\alpha_m}$. In
particular, if a local cutoff $\chi_m$ is one near that atom and
vanishes on the remaining support, then for a
bounded length weight $H$,
\begin{equation}
 H(\ell_m)=
 \frac{\langle\sigma^X_{\alpha_m},H(s)\chi_m(s)\rangle}
      {\langle\sigma^X_{\alpha_m},\chi_m(s)\rangle}.
 \label{eq:polearccoefficientweight}
\end{equation}
For arbitrary bounded Borel $H$, the same statement is interpreted
using the already recovered first-atom location. No assertion that
first-atom extraction is a bounded linear functional is being made.

Choose an even $f\in C_c^2(\RR^2)$ supported in the unit disk, a bounded
Borel function $h:(0,\infty)\to\CC$, and parameters $R\geq1$, $L>0$.
Define
\begin{equation}
 \mathcal F_{\delta;f,h}^{R,L}(q_X)
  =\frac1{2R^2}
    \sum_{\substack{m\in\ZZ^2\ \mathrm{primitive}\\m\equiv\delta\ (2)}}
       f(gm/R)\,h\bigl(\ell_m(q_X)/L\bigr).
 \label{eq:polearcobservable}
\end{equation}
The factor $1/2$ identifies opposite directions. The function is
well defined on $Y_2$: changes of lattice basis preserve the parity
class and relabel the sum. More precisely, define the geometric formula
on the invariant conull admissible set of
Lemma~\ref{lem:primitivepolearcs} and set it to zero on its complement;
this fixes a measurable representative of the same $L^2$ class.
Averaging over the finite labeling choices
gives observables on the unlabelled stratum, with the same norm bounds.
Such an average is an average of selected families; it need not count
their union without multiplicities.

\subsection{Uniform rotation norms, including a sharp hyperbolic cutoff}
\label{subsec:polearcnorms}

Put
\[
 J=\begin{pmatrix}0&-1\\1&0\end{pmatrix},\qquad
 A=\begin{pmatrix}1&0\\0&-1\end{pmatrix},\qquad
 D_Bf(x)=(Bx)\cdot\nabla f(x).
\]
The notation $D_B$ here is a linear vector field on $\RR^2$.

\begin{lemma}[Uniform primitive-vector bound]
\label{lem:polearcprimitivebound}
For every unimodular lattice $g\ZZ^2$ and every $R>0$,
\begin{equation}
 \#\{m\in\ZZ^2:\ m\ \mathrm{primitive},\ |gm|\leq R\}
 \leq 2+2\pi R^2.
 \label{eq:polearcprimitivebound}
\end{equation}
\end{lemma}

\begin{proof}
Two nonparallel lattice vectors $u,v$ have $|\det(u,v)|\geq1$.
If both have length at most $R$, their unoriented angular separation
therefore has sine at least $R^{-2}$. For $R<1$ there is at most one
primitive direction. For $R\geq1$, placing the directions on the circle
of length $\pi$ shows that their number is at most
$\pi/\arcsin(R^{-2})\leq\pi R^2$. Each direction contributes only
the two opposite primitive vectors. The displayed bound covers both
cases without rounding conventions.
\end{proof}

Rotation of the flat charts fixes the underlying conformal surface and
the intrinsic pole-to-pole arcs. Consequently
\[
 \Theta\ell_m=0,
\]
and differentiation of~\eqref{eq:polearcobservable} under rotation
only replaces $f$ by $D_Jf$. Set
\[
 C_*=1+\pi,\qquad
 N_{f,h}=C_*\|h\|_\infty
             \bigl(\|f\|_\infty^2+\|D_Jf\|_\infty^2\bigr)^{1/2}.
\]
The primitive-vector bound proves, uniformly in $R\geq1$ and $L>0$,
\begin{equation}
 \mathcal S_1(\mathcal F_{\delta;f,h}^{R,L})\leq N_{f,h}.
 \label{eq:polearcuniformS1}
\end{equation}
This is global, including the modular cusp. Although the curve labels
can jump as a regular mark crosses a selected arc, the exceptional walls
are $SL_2$-invariant. Along almost every group orbit the labels are fixed
under the natural marking, and the preceding derivatives are legitimate.
Uniform bounds on the difference quotients give the corresponding strong
$L^2$ generator statements.

In particular, Theorem~\ref{thm:markedtorusmixing} now applies to two
observables of the form~\eqref{eq:polearcobservable}, with
\begin{equation}
 |\operatorname{Corr}_t(\mathcal F,\mathcal G)|
 \leq4\coth\pi\,(1+t)e^{-t}\,N_{f,h}N_{\widetilde f,\widetilde h},
 \label{eq:polearccorrelation}
\end{equation}
uniformly in their two flat cutoffs and two hyperbolic-length scales.
The hyperbolic weight may be the sharp cutoff
$h(s)=\mathbf1_{s\leq1}$: its discontinuity costs no rotation derivative.
If the flat weight is radial as well, the whole observable is
$K$-invariant. In that case a sharp disk cutoff is also allowed directly,
and the single-$K$-type matrix-coefficient bound gives the sharper estimate
\begin{equation}
 |\operatorname{Corr}_t(F,G)|
 \leq\Xi(a_t)\|F^\circ\|_2\|G^\circ\|_2
 \leq\frac4\pi(1+t)e^{-t}\|F^\circ\|_2\|G^\circ\|_2.
 \label{eq:polearcKfixedcorrelation}
\end{equation}
No rotational smoothing is needed for these radial tests.

\subsection{Control of the terminal derivative}
\label{subsec:polearcterminal}

For the derivative counting converse, require $h\in C^1$ and
\[
 H_0=\|h\|_\infty<\infty,\qquad
 H_1=\sup_{s>0}s|h'(s)|<\infty.
\]
Wolpert's quasiconformal length inequality gives
\begin{equation}
 |V\ell_m|\leq2\ell_m,
 \qquad |\Theta V\ell_m|\leq4\ell_m.
 \label{eq:polearclengthderivatives}
\end{equation}
For completeness, the first inequality follows by lifting a
$K$-quasiconformal map to the cyclic annular covers associated with the
curve. Their conformal moduli are $\pi/\ell_X$ and $\pi/\ell_Y$;
modulus distortion gives $K^{-1}\leq\ell_Y/\ell_X\leq K$.
The map associated with $a_t$ has dilatation $e^{2|t|}$.
For the second inequality use $\Theta\ell_m=0$ and
$[\Theta,V]=\pm2W$, where $W$ is the generator of a rotated diagonal
subgroup. The same first inequality applies to $W$.
See also~\cite[Equation~(1.1)]{ShigaHyperbolicQC} for the quasiconformal
inequality and its attribution to Wolpert.

Let $F=\mathcal F_{\delta;f,h}^{R,L}$, and abbreviate
\[
 f_0=\|f\|_\infty,\quad f_K=\|D_Jf\|_\infty,\quad
 f_A=\|D_Af\|_\infty,\quad f_{KA}=\|D_JD_Af\|_\infty.
\]
Direct differentiation and~\eqref{eq:polearclengthderivatives} give
\begin{align}
 \|VF\|_\infty
  &\leq C_*\bigl(f_AH_0+2f_0H_1\bigr),\notag\\
 \|\Theta VF\|_\infty
  &\leq C_*\bigl(f_{KA}H_0+2f_KH_1+4f_0H_1\bigr).
 \label{eq:polearcVnorms}
\end{align}
These constants are again independent of $R\geq1$ and $L>0$.
No second derivative of $h$ occurs, because rotation fixes $\ell_m$.
Thus $\mathcal S_1(VF)$, and hence $\mathcal S_1((4+V)F)$, have explicit
uniform bounds. In particular the terminal identity extends to these
observables:
\begin{equation}
 \int_1^T e^{4t}\operatorname{Corr}_t(F,(4+V)G)\,\dd t
 =e^{4T}\operatorname{Corr}_T(F,G)
   -e^4\operatorname{Corr}_1(F,G).
 \label{eq:polearcterminalidentity}
\end{equation}
Here $G$ has a smooth length weight as above; invariance of probability
measure gives $\mu(VG)=0$. This proves applicability of the integrated
counting converse to these endpoints. It does not identify an individual
curve sum with the finite affine return sum of the preceding section.
Sharp cutoffs require a separate smoothing or boundary argument when
the terminal derivative is used.

\subsection{What the bridge retains}
\label{subsec:polearcscope}

For $h\equiv1$, the observable reduces to a primitive parity Siegel
transform on the modular base and has no nonzero marked-point Fourier
mode. For a disk cutoff and fixed $g$, elementary lattice counting and
M\"obius inversion give
\[
 \frac12\#\{m\ \mathrm{primitive}:m\equiv\delta\ (2),\ |gm|\leq R\}
 =\frac{\pi}{6\zeta(2)}R^2+O_g(R\log R)
 =\frac1\pi R^2+O_g(R\log R).
\]
Indeed, only odd common divisors occur in a nonzero parity class, and
$\sum_{d\ \mathrm{odd}}\mu(d)/d^2=4/(3\zeta(2))$. Adding the hyperbolic
weight $h(\ell_m/L)$ can retain dependence on the two regular marks and
uses the simple-curve first-atom data. For $h\equiv1$ the exact mean is
\[
 \mu(\mathcal F_{\delta;f,1}^{R,L})
 =\frac1{6\zeta(2)}\int_{\RR^2}f(x)\,\dd x,
\]
by the primitive Siegel formula and the equal densities of the three
nonzero parity classes.

There is a useful fixed-surface sparsity bound. On a fixed compact
$(X,q_X)$ put
\[
 C(X,q_X)=\sup_X\frac{|q_X|^{1/2}}{\dd s_{\rm hyp}}<\infty.
\]
The selected flat geodesic minimizes flat length in its free homotopy
class, so comparison with the hyperbolic geodesic gives
\[
 |gm|=\ell_{q_X}(\alpha_m)
 \leq C(X,q_X)\ell_X(\alpha_m).
\]
Consequently the whole selected family with hyperbolic length at most
$L$ lies in the flat window $R=C(X,q_X)L$ and has size
$O_{X,q_X}(L^2)$. It is a sparse family relative to the full genus-two
$L^6$ simple-curve count. The bound is for a fixed surface and
differential; no uniform comparison constant over moduli is asserted.

Equations~\eqref{eq:polearccoefficientweight}--\eqref{eq:polearcVnorms}
are a norm-controlled bridge for this selected family. They do not
identify the full connected-simple-curve count with an affine return
sum. Nor do they prove uniform stability of extracting first atoms from
perturbed spectral or coefficient data: the shrinking-gap issue of
Section~\ref{subsec:countnearcollision} remains. The output observables
have controlled geometric norms, while the quantitative conditioning of
the nonlinear coefficient extraction is a different question.

\subsection{Affine displacement sums: an exact isometric model}
\label{subsec:affinedisplacements}

The homogeneous calculation can be tested on explicit geometric sums.
These are the relative-period Siegel--Veech transforms discussed in
\cite[Section~6.1]{AthreyaLagaceMollerRaum}.
On the marked-torus model write a point as $(g,v_1,v_2)$, with the
physical marked points understood modulo $g\ZZ^2$. For
$f\in\mathscr S(\RR^2)$ define
\[
 S_{f,j}(g,v_1,v_2)=\sum_{m\in\ZZ^2}f(gm+v_j),\qquad
 F_{f,j}=S_{f,j}-\int_{\RR^2}f(x)\,\dd x.
\]
The summands test the displacement vectors of straight segments from
the torus origin to its $j$th marked point. These are flat geometric
observables; a hyperbolic simple-curve interpretation requires a
separate choice of a lift and of curve weights.

Use the Fourier transform
$\widehat f(\xi)=\int f(x)e^{-2\pi i x\cdot\xi}\dd x$.
With $v_j=gu_j$, Poisson summation gives
\[
 F_{f,j}(g,u_1,u_2)=
 \sum_{n\in\ZZ^2\setminus\{0\}}
       \widehat f(g^{-T}n)e^{2\pi i n\cdot u_j}.
\]
Fiber Parseval followed by the usual Siegel mean formula for all
nonzero lattice vectors gives the exact identity
\begin{equation}
 \|F_{f,j}\|_{L^2(Y_2)}^2=\|f\|_{L^2(\RR^2)}^2.
 \label{eq:affinesiegelisometry}
\end{equation}
This is also the one-relative-mark case of
\cite[Proposition~6.11]{AthreyaLagaceMollerRaum}.
The level-two probability cover changes neither side: the integrand
descends to the full modular quotient. By polarization and equivariance,
\begin{equation}
 \boxed{
 \int_{Y_2}\overline{F_{f,j}(x)}F_{h,j}(a_tx)\,\dd\mu(x)
 =\int_{\RR^2}\overline{f(x)}h(a_tx)\,\dd x.
 }
 \label{eq:affinesiegelexactcorr}
\end{equation}
Indeed the Fourier-side integral is
$\int\overline{\widehat f(\xi)}\widehat h(a_{-t}\xi)\dd\xi$.
These identities extend by density to $L^2$ tests, using the isometric
definition of the centered transform when the original sum is unavailable.
For indicator tests the variance identity is also
\cite[Lemma~4]{AthreyaAffineLattices}; the mean is Lemma~3 there.
For different relative marks, the analogous centered cross-correlation
vanishes by integration in the two independent torus fibers.

For the Gaussian $f(x)=e^{-\pi|x|^2}$, the formula becomes
\begin{equation}
 \int_{Y_2}\overline{F_{f,j}}F_{f,j}\circ a_t\,\dd\mu
 =\frac1{\sqrt{(1+e^{2t})(1+e^{-2t})}}
 =\frac1{2\cosh t}.
 \label{eq:affinegaussiancorr}
\end{equation}
Thus an $e^{-t}$ correlation rate is attained by a natural smooth
observable. More generally, substitution $x_1=e^{-t}u$ in
\eqref{eq:affinesiegelexactcorr} gives, for Schwartz tests,
\[
 e^t\int\overline f(x)h(a_tx)\,\dd x
 \longrightarrow
 \left(\int_\RR\overline{f(0,y)}\,\dd y\right)
 \left(\int_\RR h(u,0)\,\dd u\right).
\]
Schwartz domination justifies the limit.

\subsection{Two marks: the rank-one divergence and its removal}
\label{subsec:affineranktwo}

The raw product $F_{f_1,1}F_{f_2,2}$ need not be in $L^2$, even for
smooth rapidly decreasing tests. Its conditional squared norm is
\[
 A_1(g)A_2(g),\qquad
 A_j(g)=\sum_{n\ne0}|\widehat f_j(g^{-T}n)|^2.
\]
For Gaussians $f_1=f_2=e^{-\pi|x|^2}$ the divergence can be seen
exactly. Write the collinear frequency pairs as $(an,bn)$, where
$n$ is primitive modulo sign and $a,b\in\ZZ\setminus\{0\}$.
Primitive Siegel unfolding gives the rank-one contribution
\[
 \frac1{2\zeta(2)}\sum_{a,b\ne0}
       \int_{\RR^2}e^{-2\pi(a^2+b^2)|x|^2}\,\dd x
 =\frac1{4\zeta(2)}\sum_{a,b\ne0}\frac1{a^2+b^2}
 =\infty.
\]
The divergence is logarithmic. More generally, if both $\int f_j$ are
nonzero, Fourier continuity gives $A_j(g)\gg\varepsilon^{-1}$ when
the dual lattice has a primitive vector of length $\varepsilon\ll1$;
the cusp measure is comparable to $\varepsilon\dd\varepsilon$.
Centering has removed zero frequencies but leaves precisely this
collinear-frequency obstruction.

There is an explicit $L^2$ replacement. Retain only the full-rank
frequency matrices:
\begin{equation}
 H_{f_1,f_2}^{(2)}(g,u_1,u_2)
 =\sum_{\substack{n_1,n_2\in\ZZ^2\\\det(n_1,n_2)\ne0}}
 \widehat f_1(g^{-T}n_1)\widehat f_2(g^{-T}n_2)
 e^{2\pi i(n_1\cdot u_1+n_2\cdot u_2)}.
 \label{eq:affineranktwoobservable}
\end{equation}
The series is absolutely convergent on each fixed torus fiber.
The rank restriction is invariant under the modular transition and
under the $SL_2$ action. It is a fiber Fourier projection, not a
positivity-preserving deletion of physical segments. It is defined
fiberwise before global integrability has been proved; no orthogonal
projection on the global $L^2$ space is being applied to the raw product.

Here is a norm computation which proves global integrability. Put
$a_j=|\widehat f_j|^2$ and let $\sigma_1(k)=\sum_{d\mid k}d$.
Hermite normal form gives $\sigma_1(|k|)$ full-modular orbits of
integer matrices of determinant $k\ne0$, each with trivial stabilizer.
Thus Parseval and unfolding give
\begin{equation}
 \|H_{f_1,f_2}^{(2)}\|_2^2
 =\sum_{k\ne0}\sigma_1(|k|)I_k(a_1,a_2),
 \label{eq:affineranktwonorm}
\end{equation}
where, with probability-covolume Haar normalization,
\[
 I_k(a_1,a_2)
 =\frac1{|k|\zeta(2)}
  \int_{\RR^2\setminus\{0\}}\int_\RR
  a_1(x)a_2\left(\frac{kJx}{|x|^2}+sx\right)
  \,\dd s\,\dd x,
 \qquad J(x_1,x_2)=(-x_2,x_1).
\]
For completeness the Haar-coordinate formula behind this expression is
\[
 g=\left[x,\frac{Jx}{|x|^2}+sx\right],\qquad
 \dd g=\frac1{\zeta(2)}\,\dd x\,\dd s.
\]
The normalization follows by unfolding the primitive first-column
Siegel transform: the stabilizer shifts $s$ by integers, and its mean
is $\zeta(2)^{-1}\int f$. At level two each full-rank modular orbit
splits into six orbits while Haar measure is divided by six, leaving
\eqref{eq:affineranktwonorm} unchanged.

An explicit Schwartz bound is convenient. Define
\[
 A(f)=\sup_{\xi\in\RR^2}
       (1+|\xi|^2)^{3/2}|\widehat f(\xi)|.
\]
For $k>0$ the orbit integral is bounded by
\[
 I_{\pm k}\leq
 A(f_1)^2A(f_2)^2
 \int_G(1+k\operatorname{tr}(g^Tg))^{-3}\,\dd g
 =\frac{\pi^2 A(f_1)^2A(f_2)^2}
        {2\zeta(2)k(1+2k)^2}.
\]
One way to verify the last integral is the Gaussian identity
\[
 \int_G e^{-\alpha k\operatorname{tr}(g^Tg)}\,\dd g
 =\frac{\pi^2}{\alpha k\zeta(2)}e^{-2\alpha k},
 \qquad \alpha,k>0,
\]
followed by the Laplace representation of $(1+u)^{-3}$.
Since $\sum_{k\geq1}\sigma_1(k)k^{-3}=\zeta(3)\zeta(2)$,
\begin{equation}
 \boxed{
 \|H_{f_1,f_2}^{(2)}\|_2
 \leq\frac\pi2\sqrt{\zeta(3)}\,A(f_1)A(f_2).
 }
 \label{eq:affineranktwoschwartz}
\end{equation}

Let $Df(x)=(Jx)\cdot\nabla f(x)$ and set $B(f)=A(f)+A(Df)$.
Rotation commutes with the rank projection, so
\[
 \Theta H_{f_1,f_2}^{(2)}
 =H_{Df_1,f_2}^{(2)}+H_{f_1,Df_2}^{(2)}.
\]
Therefore
\begin{equation}
 \mathcal S_1(H_{f_1,f_2}^{(2)})
 \leq\frac\pi2\sqrt{\zeta(3)}\,B(f_1)B(f_2).
 \label{eq:affineranktwoKnorm}
\end{equation}
The preceding mixing theorem now applies, with constants controlled
explicitly by the four test functions. In particular,
\[
 \left|\left\langle H_{f_1,f_2}^{(2)},
       H_{h_1,h_2}^{(2)}\circ a_t\right\rangle\right|
 \leq\pi^2\zeta(3)\coth\pi\,(1+t)e^{-t}
       B(f_1)B(f_2)B(h_1)B(h_2).
\]
The mean is zero because every retained fiber frequency is nonzero.

\subsection{An exact Gaussian two-mark correlation}
\label{subsec:affineranktwogaussian}

Take $f_1=f_2=f=e^{-\pi|x|^2}$ and write $H^{(2)}=H_{f,f}^{(2)}$.
The determinant-orbit calculation can be done at every time:
\begin{equation}
 \boxed{
 \left\langle H^{(2)},H^{(2)}\circ a_t\right\rangle
 =\frac{\pi}{\zeta(2)\cosh t}
       \sum_{k=1}^\infty\frac{\sigma_1(k)}k
                         e^{-4\pi k\cosh t}.
 }
 \label{eq:affineranktwogaussiancorr}
\end{equation}
To see this, the product of the Fourier coefficients is
\[
 \exp\bigl[-\pi\operatorname{tr}(Z^TC_tZ)\bigr],
 \qquad C_t=I+a_{-2t},\qquad \det Z=k.
\]
Since $\sqrt{\det C_t}=2\cosh t$, an $SL_2$ change of variables
reduces this to the radial Gaussian integral with
$\alpha=2\pi\cosh t$. Summing the positive and negative determinant
orbits gives~\eqref{eq:affineranktwogaussiancorr}. At $t=0$ this gives
the exact squared norm. The identity
\[
 \sum_{k\geq1}\frac{\sigma_1(k)}k q^k
 =-\sum_{n\geq1}\log(1-q^n),\qquad 0<q<1,
\]
also expresses the answer by the ordinary partition Euler product.
The resulting correlation decays faster than exponentially in $t$.
This strong decay is a feature of these particular Gaussian tests;
it is consistent with the $e^{-t}$ one-mark example and the uniform
tempered bound.
For a finite evaluation, put $q=e^{-4\pi\cosh t}$. The elementary bound
$\sigma_1(k)/k\leq k$ gives the explicit truncation estimate
\[
 0\leq\sum_{k>N}\frac{\sigma_1(k)}kq^k
 \leq\frac{q^{N+1}((N+1)-Nq)}{(1-q)^2}.
\]
This controls truncation of the analytic series. Floating-point
evaluation of this bound is not an interval enclosure of the total
computational error, and checking finitely many times is not a proof
of the general mixing theorem.

\subsection{Sharp boundaries and the rotation norm}
\label{subsec:affinesharprotation}

The exact isometry also displays a regularity cost of sharp cutoffs.
For a bounded convex body $B\subset\RR^2$,
\[
 \|F_{\mathbf1_B,j}\circ k_\theta-F_{\mathbf1_B,j}\|_2^2
 =\operatorname{area}(B\mathbin\triangle k_\theta B).
\]
The first variation of area under rotation, interpreted using the
almost-everywhere outer normal of a convex boundary, gives
\[
 \operatorname{area}(B\mathbin\triangle k_\theta B)
 =|\theta|\int_{\partial B}|Jx\cdot n(x)|\,\dd s+o(|\theta|).
\]
Unless $B$ is rotationally invariant, this rules out a finite
$L^2$ rotation derivative. Equivalently, the sharp transform has
angular Sobolev regularity of order $s<1/2$, but not $s\geq1/2$.
The direct summation of tempered $K$-type estimates requires an
order greater than $1/2$; the earlier theorem uses order one.

This is a generic boundary-regularity phenomenon. It is not by
itself the more specific arithmetic obstruction from the unusually
flat McShane--Rivin norm ball. It explains why a smooth correlation
bound cannot be applied to a sharp norm-ball count while leaving
the observable norm unchanged. Here the body $B$ is held fixed in
physical Euclidean coordinates. A cutoff depending only on the
hyperbolic length $\ell_X(\alpha)$ is unchanged when the quadratic
differential is rotated over the same conformal surface $X$, and
does not incur this particular rotation loss. These two kinds of
cutoffs must be distinguished. The fixed-surface primitive-lattice
count, its M\"obius inversion, and its unsmoothing error remain
distinct from the affine mean-square identities proved here.

\paragraph{Finite computational companion.}
The companion \path{experiments/counting_mixing/} checks exact finite
algebra and selected one- and two-mark Gaussian formulas at $60$ and
$90$ decimal working digits, including rank-two series/orbit comparisons.
Analytic truncation bounds do not certify quadrature or roundoff.
These fixtures prove neither general mixing nor the nonsquare genus-two
model hypotheses; finite collinear growth does not prove divergence.
Its README links results, review and generation records.

